% ============================================================================
\section{Introducción}
\label{sec:intro}
% ============================================================================

Cada año los colegios entregan a las familias una lista de útiles: un documento
escrito por materias, con redacciones libres, cantidades que a veces están en
una columna y a veces dentro del texto, marcas sugeridas, instrucciones del
tipo «el disponible en casa» y libros de edición cerrada. Convertir ese
documento en una compra exige resolver tres problemas de naturaleza distinta,
que se encadenan:

\begin{enumerate}
  \item \textbf{Leer.} Pasar de un archivo (PDF o imagen), cuya estructura es
        visual ---tablas, columnas, filas que ocupan dos líneas---, a una
        secuencia de renglones, cada uno con su descripción y su cantidad.
        Además hay que decidir si el archivo corresponde a la lista de
        \emph{un} grado.
  \item \textbf{Entender.} Decidir qué pide cada renglón en términos de un
        catálogo propio y curado: a qué \emph{concepto} corresponde
        («cuaderno», «témpera», «juego de geometría»), qué atributos exige
        («tamaño oficio», «cuadriculado») y si el catálogo puede cumplirlo.
        Si no puede, hay que decirlo en lugar de ofrecer otra cosa.
  \item \textbf{Elegir.} Para cada renglón que se puede comprar hay varios
        productos admisibles, con precios, niveles de marca, empaques y
        personajes distintos. El conjunto de elecciones forma un \emph{kit}, y
        lo que una familia valora de un kit es un compromiso entre objetivos
        en conflicto: precio, calidad de marca, disponibilidad y, a veces, el
        personaje favorito de un hijo.
\end{enumerate}

Este documento describe los algoritmos con los que el sistema resuelve los
tres problemas. Cubre únicamente el \emph{núcleo algorítmico} y deja por fuera
los mecanismos de plataforma (pagos, colas de trabajo, identidad, despliegue).
El objetivo es que un lector con formación en análisis de algoritmos pueda
reconstruir cada decisión a partir de la definición formal del problema, y
juzgar si es correcta y cuánto cuesta.

\subsection{Contribuciones del documento}

\begin{itemize}
  \item Un \textbf{modelo formal único} (\cref{sec:modelo}) que fija la
        notación de todas las etapas y define el sistema como una composición
        de funciones con estados de error explícitos.
  \item Para cada etapa, la \textbf{definición del problema}, el
        \textbf{pseudocódigo}, el \textbf{análisis de complejidad} y las
        \textbf{propiedades demostrables} que la implementación cumple
        ---por ejemplo, que la percepción nunca devuelve una lista truncada, que
        la fusión léxico-semántica nunca puntúa por debajo del canal léxico o
        que consolidar la compra entre hermanos nunca aumenta el costo.
  \item Un \textbf{análisis estructural del problema de optimización}
        (\cref{sec:opt-estructura}). Bajo la formulación vigente, los cuatro
        objetivos son aditivos por línea, los extremos del frente se resuelven
        de forma voraz y exacta, y el frente completo se puede calcular en
        tiempo polinomial mediante programación dinámica. Este resultado
        explica por qué el NSGA-II iguala al óptimo en los extremos y por qué
        su aporte real está en el compromiso con presupuesto.
  \item Una \textbf{evaluación empírica} sobre 174 listas oficiales reales
        (\cref{sec:evaluacion}), con un óptimo exacto como referencia y líneas
        base que incluyen un modelo de decisión de propósito general.
  \item Un \textbf{registro de la evolución} de cada algoritmo
        (\cref{sec:evolucion} y los recuadros «Evolución» de cada sección):
        qué versión existió, qué defecto medido la hizo cambiar y qué efecto
        tuvo el cambio.
\end{itemize}

\subsection{Principios de diseño transversales}
\label{sec:principios}

Algunas decisiones se repiten en todas las etapas. Conviene enunciarlas antes
porque explican la forma de los algoritmos, y no solo su contenido.

\begin{description}[style=nextline]
  \item[Abstenerse es una respuesta válida.] En todas las etapas existe una
        salida que significa «no lo sé» o «no lo tengo», y se prefiere a una
        respuesta plausible pero incorrecta. La percepción rechaza el archivo
        completo antes que entregar una lista a medias. El emparejador se
        abstiene si ningún producto declara el atributo que el colegio exige.
        La corrección ortográfica solo corrige cuando hay \emph{una} candidata.
        La razón es la asimetría de costos: un kit incompleto que se presenta
        como completo, o un producto equivocado en el carrito, solo lo descubre
        la familia cuando recibe el pedido.
  \item[Separar decidir de presentar.] La aritmética de unidades, las líneas
        informativas que no se cobran y la consolidación son independientes del
        algoritmo que selecciona productos. Así el optimizador se puede
        reemplazar sin tocar lo que el acudiente ve.
  \item[Determinismo y reproducibilidad.] Con la misma entrada y la misma
        semilla, el resultado es idéntico byte a byte: los empates se rompen por
        claves totales (por ejemplo, el SKU) y el NSGA-II termina por número de
        generaciones y no por convergencia.
  \item[Lo curado manda sobre lo inferido.] Cuando un dato revisado por una
        persona está disponible (el concepto asignado a un renglón, el tipo de
        renglón declarado en la lista oficial), se usa en lugar de volver a
        inferirlo del texto.
  \item[Degradar en lugar de fallar, pero nunca en silencio.] Si el canal
        semántico no está disponible, la resolución continúa solo con el canal
        léxico y lo registra. Si la percepción no puede leer una página, el
        trabajo falla con un diagnóstico y se puede reintentar. No se
        fabrica una respuesta.
\end{description}

\subsection{Organización}

La \cref{sec:modelo} establece el modelo y la notación. Las
\cref{sec:percepcion,sec:normalizacion,sec:resolucion} siguen el flujo de
datos: percepción, normalización y resolución. La \cref{sec:catalogo} describe
los algoritmos que mantienen las propiedades del catálogo de las que dependen la
resolución y la optimización. La \cref{sec:optimizacion}, la más extensa,
desarrolla la optimización multiobjetivo. La \cref{sec:evaluacion} presenta la
evaluación empírica, la \cref{sec:evolucion} consolida la evolución y la
\cref{sec:discusion} discute limitaciones y trabajo futuro. Los apéndices
reúnen todos los parámetros numéricos con su justificación y la correspondencia
entre cada algoritmo y el módulo que lo implementa.
